\section{从 GPT-3 到 InstructGPT：预训练之后还缺什么}

前面的课程已经解释了怎样把算力、数据和 Transformer 组合成 base model；本讲的问题则是：一个会续写互联网文本的模型，为什么还不能直接充当 assistant？预训练最大化的是下一个 token 的似然，而产品需要的是按用户意图回答、控制风格、调用工具并在危险请求上拒绝。两种目标相关，却并不等价。

\slidepair{slide-02.jpg}{slide-03.jpg}{从课程中的 GPT-3 起点，到 instruction following 所展示的强行为控制。}{2--3}

Slides 2--3 先把“能力”和“控制”拆开。左页的 GPT-3 已经具有广泛知识与生成能力，右页的例子却显示 instruction-tuned model 能按自然语言要求改变角色、格式和推理路径。读图时应先比较的不是答案是否更长，而是同一组参数能否被 prompt 稳定地路由到目标行为；这正是 post-training 要放大的能力。

\begin{knowledgebox}{读图：instruction following 为什么是一种 remarkable control}
Base model 只需复现训练分布中常见的文本延续；assistant 则要把 prompt 当作控制信号，在大量潜在续写中选择符合任务、语气和边界的一小部分。Post-training 往往没有凭空创造这些行为，而是把预训练中已经存在但难以调用的行为变得可寻址、可重复。
\end{knowledgebox}

\slidepair{slide-04.jpg}{slide-05.jpg}{本讲的三个核心问题，以及理解现代 post-training 时必须保留的证据边界。}{4--5}

Slide 4 把课程问题压缩成数据、优化与规模三个维度：收集什么行为样本，怎样利用这些样本，以及这种控制是否也需要大规模训练。Slide 5 随即加入方法论警告：现代 frontier recipe 的公开披露远少于预训练，早期 InstructGPT、Stiennon summarization 与 Anthropic HH 反而留下了更完整的标注规范。因而本讲给出的不是一份“唯一正确配方”，而是由公开证据支持的机制框架。

\begin{warningbox}{老师强调：post-training 仍然很 artisanal}
课堂在开头把 post-training 称为带有手工作坊性质的过程：数据清洗、rubric、候选采样、训练阶段拼接和回归门槛都会改变结果。看到某个模型报告只写“做了 SFT 和 RLHF”时，不能据此认为两家的训练过程可比；未披露的数据与工程选择往往才是主要变量。
\end{warningbox}

\subsection{SFT 与 RLHF 在流水线中的位置}

有了证据边界，下一步才是建立课程地图。标准流程先用 supervised fine-tuning（SFT，监督微调）模仿 demonstrations，再用 reinforcement learning from human feedback（RLHF，人类反馈强化学习）或直接偏好优化，让模型在多个可行回答中更偏向人类选择。二者都使用行为数据，但监督对象不同。

\slidepair{slide-06.jpg}{slide-07.jpg}{标准的 SFT→RLHF 路线，以及 SFT 必须同时回答的数据与方法问题。}{6--7}

Slide 6 的 InstructGPT 流程图是本讲的主轴：预训练模型先变成能遵循指令的 SFT policy，再进入偏好优化。Slide 7 把注意力放回 SFT 的两个 ingredients。训练方法通常只是 token-level maximum likelihood；真正困难的是 demonstration 覆盖什么任务、怎样写答案、是否包含工具轨迹、安全拒绝与长尾领域，以及这些分布会把模型推向何种风格。

\begin{importantbox}{核心分解：模仿与选择}
给定 prompt $x$，SFT 使用目标回答 $y$，最大化 $\log p_\theta(y\mid x)$，回答“应该模仿什么”；preference optimization 使用同一 prompt 下的 $y^+$ 与 $y^-$，回答“两个可行回答中更偏好哪一个”。SFT 建立可用行为的支持集，RLHF/DPO 在支持集中重新分配概率。
\end{importantbox}

\subsection{先看数据里面到底有什么}

仅说“instruction data”会掩盖最重要的差异。Slide 8 要求读者先打开数据集本身：早期 benchmark conversion、Self-Instruct、真实聊天、专家编写、synthetic traces 与 agent tool calls 看起来都像 prompt-response pair，却分别训练模型复述 benchmark、模仿聊天、展示知识或遵循交互协议。后文因此从真实样本而不是数据集名称出发。

\singleslide{slide-08.jpg}{进入 SFT data 前的检查清单：内容构成与高性能数据的关键变量。}{8}

这页没有给出结论，而是规定了分析顺序：先抽样查看具体记录，再讨论规模、style、knowledge、safety 与 tool use。若跳过内容审计，只按“多少条 instruction”比较数据集，就会把一条简单分类样本和一条经过验证的多轮工具轨迹当作同等训练信号。

\begin{table}[H]
\centering
\begin{tabular}{L{0.17\textwidth}L{0.27\textwidth}L{0.46\textwidth}}
\toprule
对象 & 典型输入信号 & 在控制链中的作用 \\
\midrule
SFT & demonstrations $(x,y)$ & 建立回答格式、工具协议、安全边界与任务覆盖 \\
Preference data & comparisons $(x,y^+,y^-)$ & 测量同一 prompt 下的相对偏好 \\
Reward model / judge & preference labels & 把候选回答映射为可优化的分数或排序 \\
Reference policy & 通常为 SFT checkpoint & 约束 policy 不要为追逐 reward 而过度漂移 \\
PPO & 在线 rollouts、reward、value & 主动采样并直接提高估计 reward \\
DPO & 静态 preference pairs & 离线提高 preferred response 的相对概率 \\
Mid-training & 大规模 mixture & 在短 SFT 前重塑模型的能力与数据分布 \\
\bottomrule
\end{tabular}
\caption{本讲后续各阶段的数据接口与责任边界。}
\end{table}

\subsection{本章小结}

Post-training 不是单个 loss，而是一条从 demonstrations、comparisons 到 policy update 的控制链。预训练提供能力，SFT 使行为可调用，偏好优化重新排序候选行为；任何一环的数据定义含糊，都会让最终模型的变化难以归因。

